(1) 有效势 \(U_{\rm eff}=L^2/(2mr^2)+ar^p\)，圆轨道由 \(U'_{\rm eff}=0\) 得 \(L^2/m=ap r_c^{p+2}\)，故 \(r_c=[L^2/(map)]^{1/(p+2)}\)。

(2) 对平衡点作二阶展开，\(U''_{\rm eff}(r_c)=3L^2/(mr_c^4)+ap(p-1)r_c^{p-2}=ap(p+2)r_c^{p-2}>0\)。于是 \(\omega_r^2=ap(p+2)r_c^{p-2}/m\)。角频率 \(\omega_\phi=L/(mr_c^2)\)，由圆轨道条件得 \(\omega_\phi^2=ap r_c^{p-2}/m\)，所以 \(\omega_r/\omega_\phi=\sqrt{p+2}\)。

(3) 一次径向往返周期 \(T_r=2\pi/\omega_r\)，其间极角增加 \(\Delta\phi=\omega_\phi T_r=2\pi/\sqrt{p+2}\)。相对于闭合 Kepler 轨道的一圈进动可以写 \(\Delta\phi-2\pi\)，但不同中心势的参考闭合角须明确；本式是直接可观测的近心点间隔。

(4) \(E_c=L^2/(2mr_c^2)+ar_c^p=ap r_c^p/2+ar_c^p=a(1+p/2)r_c^p\)。当 \(p=2\)，\(\omega_r=2\omega_\phi,\Delta\phi=\pi\)，近心点在谐振子椭圆轨道的两个相反顶点之间交替，结果一致。
